\section{Lean formalization architecture}
\label{sec:lean}

The supplied Lean files are organized so that useful results can be checked
incrementally.  The small structural kernel targets declarations already
present in the repository; the remaining files state dependency-ordered
blueprints.  This section gives a concrete path from those blueprints to a
kernel-checked compiler.

\subsection{Module graph}

\begin{center}
\begin{tikzpicture}[
  node distance=8mm and 10mm,
  box/.style={draw,rounded corners,align=center,fill=darcgray,inner sep=5pt,font=\small},
  arr/.style={-{Latex[length=2mm]},thick}
]
\node[box] (up) {ConciseIPRA.PR\\IPRADerivations.Core};
\node[box,below left=of up] (struct) {DARC.StructuralKernel};
\node[box,below right=of up] (demand) {DARC.Demand};
\node[box,below=of struct] (bool) {DARC.BoundedReflection};
\node[box,below=of bool] (trace) {DARC.Trace};
\node[box,right=of trace] (rel) {DARC.Relational};
\node[box,below right=10mm and -8mm of trace] (comp) {DARC.Compiler};
\node[box,below=of comp] (bridge) {IPRABridge export checks};
\draw[arr] (up) -- (struct);
\draw[arr] (up) -- (demand);
\draw[arr] (struct) -- (bool);
\draw[arr] (bool) -- (trace);
\draw[arr] (bool) -- (rel);
\draw[arr] (demand) -- (comp);
\draw[arr] (trace) -- (comp);
\draw[arr] (rel) -- (comp);
\draw[arr] (struct) -- (comp);
\draw[arr] (comp) -- (bridge);
\end{tikzpicture}
\end{center}

The demand module is metatheory: it proves facts about finite plans and returns
resource states.  The structural, reflection, trace, relational, and compiler
modules construct values of the ordinary \code{Prf PRAAx} type.

\subsection{Structural kernel}

The file \pathcode{lean/DARC/StructuralKernel.lean} contains no
\code{sorry}, \code{admit}, or new \code{axiom}.  It defines:
\begin{itemize}
  \item \code{OpenContract} and \code{openContract\_iff};
  \item context replacement by \code{Prf.trans};
  \item free-variable specialization by \code{Prf.subst};
  \item a scoped certificate structure and sealing by \code{Prf.exE};
  \item a wrapper around \code{Search.search\_found} for named export;
  \item a quantifier-free loop structure and runner using \code{IPRA.ind};
  \item reuse of closed derivations under arbitrary hypotheses.
\end{itemize}
The first formalization step is to place this file in the current Lake
workspace and resolve any namespace or import drift.  Since every proof is a
direct call to a checked API, this step should be small and gives immediate
value even before the full compiler exists.

A representative declaration is:
\begin{lstlisting}[language={}]
def OpenContract (Gamma : List (Fm PTm n)) (A : Fm PTm n) : Prop :=
  forall Delta, (forall C in Gamma, Delta |- [PRAAx] C) ->
    Delta |- [PRAAx] A

theorem openContract_iff :
    OpenContract Gamma A <-> Gamma |- [PRAAx] A := ...
\end{lstlisting}
The actual file uses Lean notation and scoped variables.

\subsection{Demand formalization}

The metalevel demand module should avoid unnecessary abstract lattice
machinery.  Define a shape by capacities \(c_i\), a demand as a dependent
vector of \code{Fin (c i + 1)}, coordinatewise order, join, and rank.  Prove:
\begin{enumerate}
  \item join is an upper bound;
  \item strict coordinatewise increase raises rank;
  \item rank is bounded by the sum of capacities;
  \item inflationary iteration stabilizes within that sum;
  \item monotonicity implies leastness.
\end{enumerate}
For executable plans, a threshold rule can be a finite map from trigger
coordinates to levels and from required coordinates to levels.  A generic
lemma proves its transfer monotone.  More complicated modules may provide a
custom transfer and a monotonicity theorem.

The blueprint \pathcode{lean/DARC/DemandBlueprint.lean} records the intended
signatures.  Its proof milestones are metatheoretic and do not add object
axioms.

\subsection{Bounded reflection}

A robust representation should store atom evidence in the syntax.  One option
is a dependent structure:
\begin{lstlisting}[language={}]
structure ReflectedAtom (n : Nat) where
  bit      : PR n
  formula  : Fm PTm n
  range    : [] |- [PRAAx] bitTerm = 0 or bitTerm = 1
  reflect  : [] |- [PRAAx] (bitTerm = 1 <-> formula)
\end{lstlisting}
The fixed formula syntax then contains \code{ReflectedAtom}, Boolean
connectives, and bounded binders.  The generated \code{bit} and
\code{denote} functions recurse in Lean.  Reflection is proved by induction on
the syntax.  The bounded binder cases should reuse existing finite sums and
products rather than inventing a second arithmetic library.

Binder substitution is the main engineering risk.  The implementation should
first establish reusable lemmas for shifting the generated body symbol and
instantiating it at the bounded index.  Small tests with one free parameter
under two nested binders should precede the strategy compiler.

\subsection{Trace formalization}

A trace specification stores a finite state arity, initial terms, step terms,
and a quantifier-free invariant.  Generate state functions using existing
primitive recursion and pairing/simultaneous recursion.  Prove zero and
successor equations at arbitrary terms.  Then expose a theorem whose inputs
are exactly the base and one-step invariant derivations and whose output is the
open invariant.

Do not encode the invariant as an arbitrary syntax code evaluated at runtime.
It is a fixed \code{Fm PTm} value, and the theorem is generated for it.
Bounded assertions inside the informal invariant must already have been
replaced by reflected bits.

\subsection{Relational formalization}

For a first implementation, use a uniform numerical bound at each layer and a
fixed primitive-recursive edge bit.  Define viability by primitive recursion on
remaining depth, using a finite \code{any} fold.  Prove the one-step theorem:
\[
  V_h(x)=1\to
  \exists y<B_{h+1}\,[E_h(x,y)=1\wedge V_{h+1}(y)=1].
\]
Then use \code{Search.search\_found} to define the least viable successor and
prove the finite path invariant.  General finite coded sets can be added after
this numerical version works.

\subsection{Proof plan and compiler}

The global plan exists in Lean metatheory and stores local theorem values.
Because theorem types depend on node formulas, a fully generic dependent graph
can become cumbersome.  A practical staged implementation is:
\begin{enumerate}
  \item generate one Lean module per concrete plan;
  \item declare node formulas and local proof constants in that module;
  \item generate the topological schedule as ordinary Lean code using
        \code{Prf.trans};
  \item use a generic metatheoretic demand library to decide which declarations
        are emitted;
  \item only later abstract common dependent graph machinery.
\end{enumerate}
This route gives checked application pilots sooner and avoids spending months
formalizing a universal compiler representation before validating the local
interfaces.

The blueprint \pathcode{lean/DARC/CompilerBlueprint.lean} states the ideal
generic theorem, but the generated-module approach is an equally valid path to
a complete checked result.

\subsection{Axiom and dependency audit}

Every completed module should have a check file that:
\begin{enumerate}
  \item imports only the intended lower layers;
  \item prints the exact theorem types;
  \item runs \code{\#print axioms} on exported declarations;
  \item scans for \code{sorryAx};
  \item records the Lake command and repository commit;
  \item checks that the ordinary arithmetic layer does not import application
        endpoints or a second backend unexpectedly.
\end{enumerate}
A theorem is marked checked only after this audit succeeds.  Source files with
no visible placeholders but no elaboration record remain source candidates.

\subsection{Suggested integration milestones}

\begin{description}
  \item[M0.] Elaborate the structural kernel and add axiom checks.
  \item[M1.] Check demand stabilization/leastness and reproduce the JSON
        example in Lean metatheory.
  \item[M2.] Reflect conjunction, disjunction, and one bounded existential;
        then add nested alternation.
  \item[M3.] Prove the square-accumulator trace and a finite viability path as
        ordinary sequents.
  \item[M4.] Generate one concrete proof graph module and compile it by
        simultaneous cut.
  \item[M5.] Export the result through the open Bridge and compare client
        symbols.
  \item[M6.] Migrate one real packet with a nontrivial residual modulus and
        show that discharging the producer removes the boundary.
\end{description}
The detailed commands and acceptance criteria are in
\pathcode{lean/FORMALIZATION_PLAN.md}.
